% Nouveau EGA V (cote 25): archive cote 25, site page 121 (PDF page 122).
% typescript. Draft compared: none.
% AI-transcribed by Claude Opus 5.5 (claude-opus-5-5, adaptive thinking, effort high), Message Batches API, batch ega25b (msgbatch_01A7n5YuXubUPNK1zsepi8RG), 2026-11-13; not yet checked by eye.
\scanpage{25/121}
\mspage{24}

sera automatiquement \var{\deleted{de Cartier}}{}{deletion}{ega25_25_121-v0001}\var{\inserted{diviseur}}{diviseur}{insertion}{ega25_25_121-v0002} relatif pour la base $\overset{\hand{\vee}}{P}$ , comme il
résulte du fait que U contient les éléments associés au $\underline{O}_{X_s}$ , s $\in$ S,
condition qui est stable par tout changement de base et en particulier
par S'=$\overset{\hand{\vee}}{P} \longrightarrow$ S. Or il résulte aussitôt du lemme \var{\inserted{4.2.}}{4.2.}{insertion}{ega25_25_121-v0003} plus haut que D' se
prolonge en un diviseur \var{\deleted{de Cartier}}{}{deletion}{ega25_25_121-v0004} D \uline{sss} $\underline{N}'$ se prolonge en un Module
inversible, ou encore $i'_{\hand{*}}(\underline{N}')$ est inversible. \var{\deleted{Donc}}{}{deletion}{ega25_25_121-v0005} \var{\over{i}{I}}{I}{writing over}{ega25_25_121-v0006}l faudrait rédiger
mieux la fin de la démonstration en termes de cond néc et suff (sans y
aller en deux fois comme je l'ai fait), et dégager aussi \var{\inserted{d'abord}}{d'abord}{insertion}{ega25_25_121-v0007} le
\var{\margin{\unsure{voir aussi} \S 20.\illegible{1} et 20.11.}}{\placedmargin{\unsure{voir aussi} \S 20.\illegible{1} et 20.11.}}{marginal note placed}{ega25_25_121-v0008}\uline{Corollaire} \var{\inserted{3.4.}}{3.4.}{insertion}{ega25_25_121-v0009} \uline{au lemme}\hand{3.2}. Supposons g:X$\to$S plat loc de présentation finie,
soit Z une partie fermée de X telle que x $\in$ Z implique prof$_x$ $\underline{O}_{X_s} \geq 2$
\var{\inserted{soit U=X-Z et i:U$\to$X l'immersion c an.}}{soit U=X-Z et i:U$\to$X l'immersion c an.}{insertion}{ega25_25_121-v0010}
(où s = g(x)), \var{\xout{Soit $\underline{L}$ un Module loc libre sur U = X-Z, alors $\underline{L}'$}}{}{typewriter strike-out}{ega25_25_121-v0011}
\var{\xout{est isomorphe à la restriction}}{}{typewriter strike-out}{ega25_25_121-v0012} Pour tout Module loc libre \var{\inserted{de t.f}}{de t.f}{insertion}{ega25_25_121-v0013} $\underline{L}$ sur X,
considérons sa restriction $\underline{L}' = \underline{L}|U$ . Alors \var{\inserted{a)}}{a)}{insertion}{ega25_25_121-v0014} le foncteur $\underline{L} \rightsquigarrow \underline{L}'$ est
pleinement fidèle, et pour \var{\xout{\unsure{quelque} Module loc libre $\underline{L}$}}{}{typewriter strike-out}{ega25_25_121-v0015} tout $\underline{L}$, \var{\xout{xxxx}}{}{typewriter strike-out}{ega25_25_121-v0016}
l'hom canonique $\underline{L} \to i_{\hand{*}}(\underline{L}')$ est un isomorphisme. \var{\inserted{Pour que $\underline{L}$ soit de rang n, il f et s que $\underline{L}'$ le soit.}}{Pour que $\underline{L}$ soit de rang n, il f et s que $\underline{L}'$ le soit.}{insertion}{ega25_25_121-v0017} \hand{(} b) Soit $\underline{L}'$ un
Module loc libre sur X, alors $\underline{L}'$ est isomorphe à la restriction d'un
Module loc libre $\underline{L}$ sss $i_{\hand{*}}(\underline{L}')$ est loc libre. c) Supposons que $\underline{L}'$ soit
le Module inversible associé à un diviseur \var{\deleted{de Cartier}}{}{deletion}{ega25_25_121-v0018} D'\hand{/}sur U. Alors
la condition envisagée dans b) est aussi néc et suffisante pour que D'
soit la restriction d'un diviseur \var{\deleted{de Cartier}}{}{deletion}{ega25_25_121-v0019} D sur X, \var{\inserted{lequel sera alors unique,}}{lequel sera alors unique,}{insertion}{ega25_25_121-v0020}\var{\inserted{et égal à l'adhérence schématique de D' dans X.}}{et égal à l'adhérence schématique de D' dans X.}{insertion}{ega25_25_121-v0021} \hand{|} Pour que D' soit
diviseur \var{\deleted{de Cartier}}{}{deletion}{ega25_25_121-v0022} relatif par rapport à S, il f. et s. \var{\inserted{\unsure{que}}}{\unsure{que}}{insertion}{ega25_25_121-v0023} D le soit.

On utilise simplement le fait que tout $\underline{L}$ satisfait aux \var{\xout{xxx}}{}{typewriter strike-out}{ega25_25_121-v0024}
hypothèses énoncées pour $\underline{F}$ dans le lemme\hand{3.2}...

\uline{Corollaire} \var{\inserted{3.5.}}{3.5.}{insertion}{ega25_25_121-v0025} Supposons \var{\inserted{les anneaux locaux de X factoriels, (par exemple}}{les anneaux locaux de X factoriels, (par exemple}{insertion}{ega25_25_121-v0026} X régulier\hand{)} \var{\deleted{(par exemple S régulier, et X lisse sur
S)}}{}{deletion}{ega25_25_121-v0027}, Alors l'application $\emptyset \to$ \unsure{$\emptyset'$|$\emptyset$} est bijective. En particulier, si X est
un préschéma régulier loc de type fini sur un corps k, et P un fibré
projectif sur k, il y a correspondance biunivoque entre l'ensemble $\emptyset$ des

\var{\margin{121}}{\placedmargin{121}}{marginal note placed}{ega25_25_121-v0028}
